% ID: SEL-D0047-OB-E01
% ESTADO: validada. S048.
\textbf{a)} Sean $x$ e $y$ los números de envases pequeños y grandes. La capacidad impone $x+y\le1000$, el stock mínimo exige $x\ge100$ e $y\ge200$, y la relación entre demandas se traduce en $y\ge x$. El planteamiento consiste en minimizar $C=0{,}10x+0{,}20y$ euros con esas restricciones y $x,y$ enteros.

\textbf{b)} El segundo recinto es otro problema. Despejando $x$ queda $5+y/2\le x\le 10-2y$, con $y\ge0$. En $y=0$ se obtienen $(5,0)$ y $(10,0)$. Al igualar las dos cotas se obtiene $y=2$ y $x=6$. La región es el triángulo cerrado con esos tres vértices.
\begin{center}
\setlength{\unitlength}{0.9mm}
\begin{picture}(94,40)(-8,-8)
\qbezier(0,0)(42,0)(84,0)
\qbezier(0,0)(0,15)(0,30)
\put(84,1){$x$}
\put(2,28){$y$}
\qbezier(35,0)(52.5,0)(70,0)
\qbezier(35,0)(38.5,9)(42,18)
\qbezier(42,18)(56,9)(70,0)
\put(30,-6){\scriptsize $(5,0)$}
\put(64,-6){\scriptsize $(10,0)$}
\put(39,20){\scriptsize $(6,2)$}
\end{picture}
\end{center}
